\documentclass{FUCKJournalCN}

\FJtitle{请在此处填写论文标题}
\FJsubtitle{如有需要，可添加副标题}
\FJauthor{第一作者$^{1,*}$，第二作者$^{1}$，第三作者$^{2}$}
\FJaffiliation{$^{1}$单位 / 学院，大学名称，城市，邮编，国家\\
$^{2}$单位 / 学院，大学名称，城市，邮编，国家}
\FJabstract{请用简洁的学术语言概述研究问题、材料或方法、核心论点与主要结论。}
\FJkeywords{普遍性；复杂性；系统；知识}

\begin{document}

\makefjtitle

\section{引言}
说明研究背景、问题意识以及本文的核心论题。

\section{论证与方法}
\subsection{概念框架}
介绍本文采用的关键概念、材料范围和分析路径。

\subsection{分析展开}
按照清晰的小节结构组织正文论证。

\begin{equation}
\int_{0}^{1} x^2 \, dx = \frac{1}{3}
\end{equation}

\begin{figure}[htbp]
\centering
\fbox{\rule{0pt}{4cm}\rule{0.8\columnwidth}{0pt}}
\caption{请将此占位图替换为你的图像。}
\label{fig:placeholder}
\end{figure}

\begin{table}[htbp]
\centering
\caption{请将此占位表替换为你的表格。}
\begin{tabularx}{\linewidth}{>{\raggedright\arraybackslash}X>{\raggedleft\arraybackslash}X}
\toprule
项目 & 数值 \\
\midrule
示例行 & 1 \\
\bottomrule
\end{tabularx}
\end{table}

\section{结论}
总结本文贡献、边界与后续可展开的问题。

\begin{thebibliography}{99}
\bibitem{sample1}
作者. (2025). \textit{示例文献标题}. 出版社.
\end{thebibliography}

\end{document}
